%% kapa: exam=1 round=23 year=2012 period=2 form=A subject=appraisal_law subject_ko=부동산관계법규 seq=31 number=71
%% source: Q-Net 감정평가사 기출문제 (cst00309) artlSeq=5052580 file=raw/1008/5052580/2012년 제23회 감정평가사 제1차시험 기출문제(2교시 A형).pdf
%% conversion: kordoc 4.24.51 -> tools/gen1.py (auto); review=proofread-sonnet; flags=none
\begin{question}{71}
건축법령상 건축물의 구조 및 재료에 관한 설명으로 옳지 않은 것은?
\begin{choices}
  \item 건축물은 고정하중, 적재하중, 적설하중 등에 대하여 안전한 구조를 가져야 한다.
  \item 인접 대지경계선으로부터 직선거리 1.5 m에 이웃 주택의 내부가 보이는 창문을 설치하려면 차면시설을 설치하여야 한다.
  \item 방화지구 안에서 간판, 광고탑을 건축물의 지붕위에 설치하는 경우 그 높이가 3 m 미만인 경우에는 주요부를 불연재료로 하지 않아도 된다.
  \item 바닥면적의 합계가 3천,m$^2$인 공연장을 지하층에 설치하는 경우에는 각 실에 있는 자가 피난층으로 대피할 수 있도록 천장이 개방된 외부 공간을 설치하여야 한다.
  \item 건축물의 일부가 문화 및 집회시설에 해당하는 경우에는 그 부분과 다른 부분을 방화구획으로 구획하여야 한다.
\end{choices}
\end{question}
\answer{3}
